\section{Sudden and Adiabatic Limits (11/6)}

In this section we study the time evolution of a quantum state under the time-dependent Schrodinger equation 
\begin{equation}
\dket{\psi(t)}=
\mathcal T\exp\left[-\frac{\imth}{\hbar}
\int_0^t\hat H(\tau)\dif\tau\right]\dket{\psi(0)},
\qquad
\imth\hbar\frac{\dif}{\dif t}\dket{\psi(t)}
=\hat H(t)\dket{\psi(t)}.
\end{equation}
Specifically we focus on two limit cases: the extremely fast limit, described by the sudden approximation; and the extremely slow limit, described by the adiabatic evolution. 

\subsection{Sudden approximation}
In a very short time $t$
\begin{equation}
 t\ll\frac{\hbar}{\Delta E},
\end{equation}
where $\Delta E$ is the smallest level spacing between any two energy levels in the spectrum of the Hamiltonian, the unitary time evolution operator is approximately a global phase factor 
\begin{equation}
\hat U(t,0)=\mathcal T
\exp\left[-\frac{\imth}{\hbar}\int_0^t\hat H(\tau)\dif\tau\right]
\simeq e^{\imth\phi(t)}\id,
\end{equation}
and hence the state is approximately the same up to a phase factor
\begin{equation}
\boxed{\dket{\psi(t)}\simeq\dket{\psi(0)}.}
\end{equation}
For example,
\begin{equation}
\hat H(t)=\hat H_0+\delta\hat H\,\Theta(t)
\end{equation}
is a toy model that describes a sudden change from $\hat H_0$ to $\hat H_0+\delta\hat H$. In this toy model, it takes zero time for the Hamiltonian to jump from $\hat H_0$ at $t=-\epsilon$ to $\hat H$ at $t=+\epsilon$, and we have
\begin{align}
    \dket{\psi(t=-\epsilon)}\approx\dket{\psi(t=+\epsilon)}
\end{align}

\subsubsection{Sudden quench of a harmonic oscillator}
Consider
\begin{equation}
\hat H(t)=\hbar\omega(t)
\left(\hat a^\dagger\hat a+\frac12\right),
\qquad
\omega(t)=
\begin{cases}
\omega_0,&t<0,\\
\omega,&t>0.
\end{cases}
\end{equation}
Let
\begin{equation}
R_0=\sqrt{\frac{\hbar}{m\omega_0}},
\qquad
R=\sqrt{\frac{\hbar}{m\omega}}.
\end{equation}
The old annihilation operator is
\begin{align}
\hat a_0
&=\frac{\hat x}{\sqrt2R_0}
+\frac{\imth R_0\hat p}{\sqrt2\hbar}\nonumber\\
&=\frac12\left(\sqrt{\frac{\omega}{\omega_0}}
+\sqrt{\frac{\omega_0}{\omega}}\right)\hat a
+\frac12\left(\sqrt{\frac{\omega}{\omega_0}}
-\sqrt{\frac{\omega_0}{\omega}}\right)\hat a^\dagger.
\end{align}
Define the squeezing parameter $r$ as 
\begin{equation}
\frac{\omega_0}{\omega}=e^{2r}.
\end{equation}
Then
\begin{equation}
\hat a_0=\cosh r\,\hat a+\sinh r\,\hat a^\dagger.
\end{equation}
Since
\begin{equation}
\hat a_0\dket{0}_{\omega_0}=0,
\end{equation}
one has
\begin{equation}
\left(\hat a+\tanh r\,\hat a^\dagger\right)
\dket{0}_{\omega_0}=0.
\end{equation}
Thus the old ground state $\dket{0}_{\omega_0}$ of a SHO with frequency $\omega_0$ is a squeezed state of the new oscillator with frequency $\omega$.

\subsection{Schrodinger equation in the instantaneous eigenstate basis}
Define the instantaneous eigenstates $\{\dket{n,t}\}$ of a time-dependent Hamiltonian $\hat H(t)$:
\begin{equation}\label{instantaneous eigenstates}
\hat H(t)\dket{n,t}=E_n(t)\dket{n,t},~~~\expval{m,t|n,t}=\delta_{m,n}
\end{equation}
Expand
\begin{equation}
\dket{\psi(t)}=
\sum_n C_n(t)
\exp\left[-\frac{\imth}{\hbar}\int_0^tE_n(\tau)\dif\tau\right]
\dket{n,t}.
\end{equation}
Then the Schrodinger equation becomes
\begin{equation}
\dot C_m=-\sum_n
\exp\left[-\imth\int_0^t\omega_{nm}(\tau)\dif\tau\right]
C_n\dbra{m,t}\frac{\dif}{\dif t}\dket{n,t},
\end{equation}
where
\begin{equation}
\omega_{nm}(t)=\frac{E_n(t)-E_m(t)}{\hbar}.
\end{equation}
Differentiating the instantaneous eigenvalue equation (\ref{instantaneous eigenstates}) gives, for $m\neq n$,
\begin{equation}
\dbra{m,t}\frac{\dif}{\dif t}\dket{n,t}
=\frac{\dbra{m,t}\dot{\hat H}(t)\dket{n,t}}
{E_n(t)-E_m(t)}.
\end{equation}
Thus
\begin{align}
\dot C_m
=&-C_m\dbra{m,t}\frac{\dif}{\dif t}\dket{m,t}\nonumber\\
&-\sum_{n\neq m}
\exp\left[-\imth\int_0^t\omega_{nm}(\tau)\dif\tau\right]
C_n\frac{\dbra{m,t}\dot{\hat H}(t)\dket{n,t}}
{E_n(t)-E_m(t)}.
\end{align}

\subsection{Adiabatic theorem and Berry phase}
Typically the time-dependence of a Hamiltonian comes from one (or more) time-dependent parameter(s):
\begin{equation}
\hat H(t)=\hat H(\lambda(t)),
\qquad
\dot{\hat H}=\frac{\partial\hat H}{\partial\lambda}\dot\lambda.
\end{equation}
A sufficient adiabatic condition is
\begin{equation}
\boxed{
\left|\frac{\hbar\dot\lambda
\dbra{m}\partial_\lambda\hat H\dket{n}}
{(E_n-E_m)^2}\right|\ll1,
\qquad m\neq n.}
\end{equation}
Neglecting transitions between distinct instantaneous levels,
\begin{equation}
\dot C_m=-C_m\dot\lambda
\dbra m\partial_\lambda\dket m.
\end{equation}
Therefore, in the adiabatic limit we have
\begin{equation}
C_m(t)=C_m(0)
\exp\left[-\int_{\lambda(0)}^{\lambda(t)}
\dbra m\partial_\lambda\dket m\dif\lambda\right].
\end{equation}
Define the Berry connection
\begin{equation}
\boxed{\mathcal A_m(\lambda)
=-\imth\dbra m\partial_\lambda\dket m.}
\end{equation}
It is real because
\begin{equation}
\partial_\lambda\expval m m
=\expval{\partial_\lambda m}{m}
+\dbra{m}\partial_\lambda\dket{m}=0.
\end{equation}
The geometric phase is
\begin{equation}
\boxed{e^{\imth\gamma_m}
=\exp\left[-\imth\int\mathcal A_m(\lambda)\dif\lambda\right].}
\end{equation}
For example, for a qubit in a slowly varying magnetic field,
\begin{equation}
\hat H=-\hbar\omega\,\hat{\bm n}(t)\cdot\bm\sigma,
\qquad |\hat{\bm n}|=1,
\end{equation}
the Berry phase around a closed loop is proportional to the solid angle
subtended on the Bloch sphere.

